\section{Circuit Complexity --- DAGCircuit Sat}
%%% generated by the blueprint pipeline from TCSlib.Complexity.CircuitComplexity.DAGCircuitSat
\subsection*{Circuit Complexity --- DAGCircuit Sat: Introduction}

Declarations of the Lean module \texttt{TCSlib.Complexity.CircuitComplexity.DAGCircuitSat}.

\subsection*{Circuit Complexity --- DAGCircuit Sat: Results}

\begin{definition}[Satisfiable circuit]
\lean{BoolCircuit.DAGCircuit.Satisfiable}
\label{BoolCircuit.DAGCircuit.Satisfiable}
\leanok
\uses{BoolCircuit.DAGCircuit, BoolCircuit.DAGCircuit.eval}
A directed acyclic Boolean circuit on $n$ input bits is satisfiable if there is an
assignment of Boolean values to its $n$ inputs for which its output is true.
\end{definition}

\begin{definition}[Satisfiable fan-in-two circuits]
\lean{BoolCircuit.dagCktSat}
\label{BoolCircuit.dagCktSat}
\leanok
\uses{BoolCircuit.DAGCircuit, BoolCircuit.DAGCircuit.IsFaninTwo, BoolCircuit.DAGCircuit.Satisfiable}
The set consists of all pairs $(n,C)$, where $n$ is a natural number and $C$ is a
directed acyclic Boolean circuit on $n$ inputs, such that every gate of $C$ has at most
two incoming arguments and some assignment to the inputs makes its output true.
\end{definition}

\begin{theorem}[DAG Circuit Satisfiability Membership]
\lean{BoolCircuit.mem_dagCktSat_iff}
\label{BoolCircuit.mem_dagCktSat_iff}
\leanok
\uses{BoolCircuit.DAGCircuit, BoolCircuit.DAGCircuit.IsFaninTwo, BoolCircuit.DAGCircuit.Satisfiable, BoolCircuit.dagCktSat}
\difficulty{1}
A variable-count and Boolean DAG circuit pair belongs to the set of fan-in-two
satisfiable DAG circuits if and only if its circuit is fan-in-two and satisfiable.
\end{theorem}

\begin{definition}[Gate clauses for Tseitin encoding]
\lean{BoolCircuit.DAGGate.tseitin}
\label{BoolCircuit.DAGGate.tseitin}
\leanok
\uses{BoolCircuit.DAGGate}
For a gate $g$ and a natural-number variable $z$, its Tseitin encoding is a conjunctive
normal form over natural-number variables defined by the following clauses, with a
literal $(v,	ext{true})$ denoting $v$ and $(v,	ext{false})$ denoting its negation. A
zero-argument AND gate gives the unit clause $(z,	ext{true})$, while a zero-argument OR
or NOT gate gives $(z,	ext{false})$. For one argument $j$, AND and OR give the clauses
$(z,	ext{false})\vee(j,	ext{true})$ and $(z,	ext{true})\vee(j,	ext{false})$; NOT gives
$(z,	ext{true})\vee(j,	ext{true})$ and $(z,	ext{false})\vee(j,	ext{false})$. For two
arguments $j,k$, AND gives the four clauses
$(z,	ext{false})\vee(j,\text{false})\vee(k,\text{true})$,
$(z,\text{false})\vee(j,\text{true})\vee(k,\text{false})$,
$(z,\text{false})\vee(j,\text{true})\vee(k,\text{true})$, and
$(z,\text{true})\vee(j,\text{false})\vee(k,\text{false})$; OR gives
$(z,\text{false})\vee(j,\text{true})\vee(k,\text{true})$,
$(z,\text{true})\vee(j,\text{true})\vee(k,\text{false})$,
$(z,\text{true})\vee(j,\text{false})\vee(k,\text{true})$, and
$(z,\text{true})\vee(j,\text{false})\vee(k,\text{false})$; NOT gives
$(z,\text{true})\vee(j,\text{true})\vee(k,\text{true})$,
$(z,\text{true})\vee(j,\text{true})\vee(k,\text{false})$,
$(z,\text{true})\vee(j,\text{false})\vee(k,\text{true})$, and
$(z,\text{false})\vee(j,\text{false})\vee(k,\text{false})$. Every other argument list
gives the empty CNF.
\end{definition}

\begin{theorem}[Tseitin encoding characterizes gate evaluation]
\lean{BoolCircuit.DAGGate.eval_tseitin}
\label{BoolCircuit.DAGGate.eval_tseitin}
\leanok
\uses{BoolCircuit.DAGGate, BoolCircuit.DAGGate.eval, BoolCircuit.DAGGate.tseitin}
\difficulty{5}
Let $g$ be a gate with at most two inputs, let $z$ be a natural-number index, and let
$a$ assign a Boolean value to every natural-number index. For a Boolean value list $v$,
suppose the value of $v$ at each input index of $g$, with out-of-range indices treated
as false, equals the value assigned there by $a$. Then the Tseitin encoding of $g$ at
$z$ is satisfied by $a$ if and only if the value assigned to $z$ equals the evaluation
of $g$ on $v$.
\end{theorem}

\begin{theorem}[Clause length in the Tseitin encoding]
\lean{BoolCircuit.DAGGate.length_le_of_mem_tseitin}
\label{BoolCircuit.DAGGate.length_le_of_mem_tseitin}
\leanok
\uses{BoolCircuit.DAGGate, BoolCircuit.DAGGate.tseitin}
\difficulty{4}
For any natural-number index $z$ and any gate $g$, every clause in the Tseitin encoding
of $g$ at $z$ has length at most $3$.
\end{theorem}

\begin{theorem}[Bound on Tseitin clause count]
\lean{BoolCircuit.DAGGate.length_tseitin_le}
\label{BoolCircuit.DAGGate.length_tseitin_le}
\leanok
\uses{BoolCircuit.DAGGate, BoolCircuit.DAGGate.tseitin}
\difficulty{4}
For every natural number $z$ and every gate $g$, the CNF formula produced by the Tseitin
construction for $g$ with output variable $z$ has at most four clauses.
\end{theorem}

\begin{theorem}[Variables in the Tseitin clauses]
\lean{BoolCircuit.DAGGate.var_mem_of_mem_tseitin}
\label{BoolCircuit.DAGGate.var_mem_of_mem_tseitin}
\leanok
\uses{BoolCircuit.DAGGate, BoolCircuit.DAGGate.tseitin}
\difficulty{4}
For every natural number $z$ and gate $g$, each literal in each clause of the Tseitin
formula for $g$ with output variable $z$ has a variable that is either $z$ or one of the
argument indices of $g$.
\end{theorem}

\begin{definition}[Encoding a list of gates]
\lean{BoolCircuit.tseitinGates}
\label{BoolCircuit.tseitinGates}
\leanok
\uses{BoolCircuit.DAGGate, BoolCircuit.DAGGate.tseitin}
Given a natural number $z$, the Tseitin encoding of a list of gates is the concatenation
of their individual gate encodings: the first gate is encoded using variable $z$, and
each subsequent gate uses the next natural-number variable in sequence.
\end{definition}

\begin{theorem}[Evaluation of a list of Tseitin constraints]
\lean{BoolCircuit.eval_tseitinGates}
\label{BoolCircuit.eval_tseitinGates}
\leanok
\uses{BoolCircuit.DAGGate, BoolCircuit.DAGGate.tseitin, BoolCircuit.tseitinGates}
\difficulty{4}
Let $a$ assign a Boolean value to each natural-numbered vertex, let $z$ be a natural
number, and let $gs$ be a finite list of gates. The conjunction of all clauses in the
Tseitin encodings of these gates, with the $i$th gate encoded at vertex $z+i$, evaluates
to true under $a$ if and only if each gate's Tseitin encoding at its corresponding
vertex evaluates to true.
\end{theorem}

\begin{theorem}[Clause occurrence in a Tseitin encoding]
\lean{BoolCircuit.exists_of_mem_tseitinGates}
\label{BoolCircuit.exists_of_mem_tseitinGates}
\leanok
\uses{BoolCircuit.DAGGate, BoolCircuit.DAGGate.tseitin, BoolCircuit.tseitinGates}
\difficulty{4}
Let $z$ be a natural number, let $gs$ be a finite list of Boolean circuit gates, and let
$c$ be a clause occurring in the concatenated Tseitin encodings of those gates, where
gate $i$ is encoded using variable index $z+i$. Then there is an index $i$ within the
list such that $c$ occurs in the Tseitin encoding of gate $i$.
\end{theorem}

\begin{theorem}[Clause count for gate Tseitin encoding]
\lean{BoolCircuit.length_tseitinGates_le}
\label{BoolCircuit.length_tseitinGates_le}
\leanok
\uses{BoolCircuit.DAGGate, BoolCircuit.DAGGate.length_tseitin_le, BoolCircuit.tseitinGates}
\difficulty{4}
For every vertex index $z\in\mathbb{N}$ and every finite list of Boolean circuit gates,
the Tseitin CNF formed by encoding the gates successively, starting at index $z$,
contains at most four clauses per gate.
\end{theorem}

\begin{definition}[Circuit CNF encoding]
\lean{BoolCircuit.DAGCircuit.toCNF}
\label{BoolCircuit.DAGCircuit.toCNF}
\leanok
\uses{BoolCircuit.DAGCircuit, BoolCircuit.tseitinGates}
For a directed acyclic Boolean circuit on $n$ inputs, its conjunctive normal form
consists of the unit clause asserting that the output variable is true, followed by the
Tseitin encodings of its gates in order, starting with variable $n$.
\end{definition}

\begin{theorem}[Width of the CNF encoding]
\lean{BoolCircuit.DAGCircuit.toCNF_widthAtMost}
\label{BoolCircuit.DAGCircuit.toCNF_widthAtMost}
\leanok
\uses{BoolCircuit.DAGCircuit, BoolCircuit.DAGCircuit.toCNF, BoolCircuit.DAGGate.length_le_of_mem_tseitin, BoolCircuit.exists_of_mem_tseitinGates, Std.Sat.CNF.WidthAtMost}
\difficulty{4}
Let $C$ be a directed acyclic Boolean circuit on $n$ inputs. Every clause in the CNF
formula obtained from $C$ contains at most three literals.
\end{theorem}

\begin{theorem}[CNF length bound by gate count]
\lean{BoolCircuit.DAGCircuit.toCNF_length_le_gates}
\label{BoolCircuit.DAGCircuit.toCNF_length_le_gates}
\leanok
\uses{BoolCircuit.DAGCircuit, BoolCircuit.DAGCircuit.toCNF, BoolCircuit.length_tseitinGates_le}
\difficulty{3}
Let $C$ be a directed acyclic Boolean circuit on $n$ inputs. The length of the CNF
formula obtained from $C$ is at most four times the number of gates in $C$, plus one.
\end{theorem}

\begin{theorem}[CNF length bound by circuit size]
\lean{BoolCircuit.DAGCircuit.toCNF_length_le}
\label{BoolCircuit.DAGCircuit.toCNF_length_le}
\leanok
\uses{BoolCircuit.DAGCircuit, BoolCircuit.DAGCircuit.size, BoolCircuit.DAGCircuit.toCNF, BoolCircuit.DAGCircuit.toCNF_length_le_gates}
\difficulty{3}
Let $C$ be a directed acyclic Boolean circuit on $n$ inputs, with size defined as $n$
plus its number of gates. The length of the CNF formula obtained from $C$ is at most
$4\,\mathrm{size}(C)+1$.
\end{theorem}

\begin{theorem}[Variable indices in the CNF encoding]
\lean{BoolCircuit.DAGCircuit.toCNF_var_lt_size}
\label{BoolCircuit.DAGCircuit.toCNF_var_lt_size}
\leanok
\uses{BoolCircuit.DAGCircuit, BoolCircuit.DAGCircuit.size, BoolCircuit.DAGCircuit.toCNF, BoolCircuit.DAGGate.var_mem_of_mem_tseitin, BoolCircuit.exists_of_mem_tseitinGates}
\difficulty{5}
Let $C$ be a directed acyclic Boolean circuit on $n$ inputs, with size equal to $n$ plus
its number of gates. Every literal in every clause of the CNF formula obtained from $C$
has variable index strictly less than $\mathrm{size}(C)$.
\end{theorem}

\begin{theorem}[Circuit values satisfy gate clauses]
\lean{BoolCircuit.DAGCircuit.eval_tseitinGates_values}
\label{BoolCircuit.DAGCircuit.eval_tseitinGates_values}
\leanok
\uses{BoolCircuit.DAGCircuit, BoolCircuit.DAGCircuit.values, BoolCircuit.DAGCircuit.values_getD_gate, BoolCircuit.DAGGate.eval_tseitin, BoolCircuit.eval_tseitinGates, BoolCircuit.tseitinGates}
\difficulty{3}
Let $C$ be a directed acyclic Boolean circuit on $n$ inputs in which each gate has at
most two arguments, and let $x\in\{0,1\}^n$. Assign each vertex the value it receives
when the circuit is evaluated on $x$. The conjunction of the Tseitin clauses for all
gates of $C$ is true under this assignment.
\end{theorem}

\begin{theorem}[Tseitin assignments determine circuit values]
\lean{BoolCircuit.DAGCircuit.values_getD_eq_of_eval_tseitinGates}
\label{BoolCircuit.DAGCircuit.values_getD_eq_of_eval_tseitinGates}
\leanok
\uses{BoolCircuit.DAGCircuit, BoolCircuit.DAGCircuit.values, BoolCircuit.DAGCircuit.values_getD_gate, BoolCircuit.DAGCircuit.values_getD_input, BoolCircuit.DAGGate.eval_tseitin, BoolCircuit.eval_tseitinGates, BoolCircuit.tseitinGates}
\difficulty{6}
Let $C$ be a Boolean directed acyclic circuit with $n$ inputs, and suppose every gate
has at most two arguments. If a Boolean assignment satisfies the conjunction of the
Tseitin clauses for all gates of $C$, then for every vertex $v<n+|C.	ext{gates}|$, the
value computed at $v$ from the input assignment agrees with the given assignment at $v$.
\end{theorem}

\begin{theorem}[Circuit and CNF satisfiability]
\lean{BoolCircuit.DAGCircuit.satisfiable_iff_toCNF_satisfiable}
\label{BoolCircuit.DAGCircuit.satisfiable_iff_toCNF_satisfiable}
\leanok
\uses{BoolCircuit.DAGCircuit, BoolCircuit.DAGCircuit.Satisfiable, BoolCircuit.DAGCircuit.eval, BoolCircuit.DAGCircuit.eval_tseitinGates_values, BoolCircuit.DAGCircuit.toCNF, BoolCircuit.DAGCircuit.values, BoolCircuit.DAGCircuit.values_getD_eq_of_eval_tseitinGates, Std.Sat.CNF.Satisfiable}
\difficulty{4}
Let $C$ be a finite acyclic Boolean circuit. Suppose every gate of $C$ has at most two
input vertices. Then $C$ has an input assignment for which its output is true if and
only if the CNF formula associated with $C$ has a satisfying assignment.
\end{theorem}

\begin{theorem}[Satisfiability for fan-in-two circuits]
\lean{BoolCircuit.DAGCircuit.satisfiable_iff_toCNF_satisfiable_of_isFaninTwo}
\label{BoolCircuit.DAGCircuit.satisfiable_iff_toCNF_satisfiable_of_isFaninTwo}
\leanok
\uses{BoolCircuit.DAGCircuit, BoolCircuit.DAGCircuit.IsFaninTwo, BoolCircuit.DAGCircuit.Satisfiable, BoolCircuit.DAGCircuit.satisfiable_iff_toCNF_satisfiable, BoolCircuit.DAGCircuit.toCNF, Std.Sat.CNF.Satisfiable}
\difficulty{1}
Let $C$ be a finite acyclic Boolean circuit. Suppose each gate has no repeated input
vertices, and every negation gate has exactly one input vertex; suppose also that every
gate has at most two input vertices. Then $C$ has an input assignment for which its
output is true if and only if the CNF formula associated with $C$ has a satisfying
assignment.
\end{theorem}

\begin{theorem}[DAG Circuit Satisfiability and Its CNF Encoding]
\lean{BoolCircuit.mem_dagCktSat_iff_toCNF}
\label{BoolCircuit.mem_dagCktSat_iff_toCNF}
\leanok
\uses{BoolCircuit.DAGCircuit, BoolCircuit.DAGCircuit.IsFaninTwo, BoolCircuit.DAGCircuit.satisfiable_iff_toCNF_satisfiable_of_isFaninTwo, BoolCircuit.DAGCircuit.toCNF, BoolCircuit.DAGCircuit.toCNF_widthAtMost, BoolCircuit.dagCktSat, BoolCircuit.mem_dagCktSat_iff, Std.Sat.CNF.Satisfiable, Std.Sat.CNF.WidthAtMost}
\difficulty{3}
Let $C$ be a Boolean DAG circuit on $n$ input variables, and suppose that $C$ is
fan-in-two. Then the pair consisting of $n$ and $C$ belongs to the set of fan-in-two
satisfiable DAG circuits if and only if the CNF formula associated to $C$ has width at
most three and is satisfiable.
\end{theorem}
